\chapter{Conduits}
\label{ch:conduits}

Wire sags. Strung across a valley that is exactly right; strung along a ceiling it is exactly wrong.
\textbf{Conduit} is the indoor answer: thin tubing that clips flat against whatever face you place
it on and runs in that face, turning in right angles.

\begin{iefork}
Not in stock Immersive Engineering. It takes the bundled-cable idea from SimpleLogic and puts it in
IE's idiom --- surface conduit, junction boxes, and ordinary IE connectors at the breakouts.
\end{iefork}

One length of conduit carries \textbf{sixteen} independent conductors, named after the sixteen dyes.

\section{Runs}

Place conduit against a wall, floor or ceiling and it clips to whatever you clicked. Lengths joined
along the same surface make a \textbf{run}.

To continue a run, click the end of it. A conduit placed against another conduit joins that
conduit's surface rather than clipping to the conduit itself, so you can lay a corridor by clicking
along the run instead of hunting for the bare wall between lengths. The first length off a junction
box does the same thing --- it looks for the surface behind it rather than clipping to the box.

\textbf{A run stays on one surface.} Conduit on a floor does not climb the wall by itself. That is
deliberate: a plane change goes through a junction box, which is a real block and what you would put
at a corner anyway.

Runs are discovered, not drawn. Laying conduit between two boxes creates the connection; pulling a
length out of the middle breaks it. There is no coil and no linking tool --- the run is the thing
you built.

\begin{mbplan}{A run seen from above, on a floor}{C = conduit, J = junction box. The run turns in
the plane of the surface it is clipped to; leaving that surface needs the box.}
\mblayer{}{J,C,C,C,C,J;.,.,.,.,.,C;.,.,.,.,.,J}
\end{mbplan}

\section{The Junction Box}

A patch panel with six faces. Each face is one of three things:

\begin{itemize}
\item Right-click with a \textbf{dye}: that face breaks out the conductor of that colour.
\item Right-click a patched face with \textbf{redstone dust}: cycles it between power, reading
      redstone, and emitting redstone.
\item Sneak and right-click empty-handed: unpatches the face.
\end{itemize}

\textbf{A patched face wears a plate in its conductor's colour}, so a box tells you how it is wired
from across the room: which faces break out, and which channel each one carries. A box with no
plates on it is a box nobody has patched.

Put an LV, MV or HV Connector against a power breakout and it picks up that conductor and nothing
else. One bundle down a corridor can feed a lighting circuit, a workshop and a furnace, each on its
own colour.

\begin{ietip}
There is no tier setting. The tier of a circuit is whichever connector you hang on its breakout,
because IE's connectors already cap throughput by tier. Mixed tiers in one bundle are fine.
\end{ietip}

A colour lives on one face at a time --- patching it somewhere new takes it off wherever it was,
because the same conductor arriving at two connectors is a short.

A box passes the whole bundle through whatever is patched, so one dropped in to turn a corner or
change surface needs no configuration at all.

\section{The Ground Feeder}

A junction box is the right way to leave a surface indoors, where it reads as hardware bolted to a
wall. Out in a field it is a steel box sitting in a lawn, and there is no surface to change
\emph{to} until the run is already underground.

The \textbf{Ground Feeder} is the other answer: a whole cube that fills a hole in a floor, a wall or
the ground, passes a bundle straight through along one axis, and \textbf{draws itself as whatever is
around it}.

\textbf{It looks at what is visible rather than at what is merely nearby.} A feeder set into a grass
field has eight grass blocks around it and nine dirt blocks underneath, so counting neighbours would
put dirt in the middle of a lawn. Instead a block scores triple if any face of it can be seen at all,
and double if it shares a face with the feeder rather than meeting it at a corner. Turf wins on a
lawn; stone wins in a cellar wall, because down there the stone is what is exposed to the shaft.

Right-click a feeder with any full block to pin that block instead. Sneak and right-click it
empty-handed to hand it back to the survey.

\begin{ietip}
It can only wear full, opaque, ordinary cubes. A feeder is a cube and always will be, so wearing a
fence would draw a fence with a solid hitbox and wearing a chest would draw a chest whose lid never
opens --- both read as a bug rather than as a disguise.
\end{ietip}

A run passes through along \textbf{one axis}, which an Engineer's Hammer turns. A feeder laid across
the grain of a run stops it exactly as a stone block would, and stops it visibly: the conduit does
not draw its arm into one it cannot enter.

\textbf{Crossing a feeder is the one place a run may change surface without a box.} The conduit on
the far side still has to have the crossing in its own plane --- the same thing a junction box asks
--- so a run can come down a north wall and continue down an east one, but it cannot continue into
conduit lying flat on the floor of the shaft.

\begin{mbplan}{A run leaving a building, seen from the side}{C = conduit on a wall, G = ground
feeder, J = junction box. The feeder fills the floor; the run crosses it and carries on below.}
\mblayer{}{J;C;G;C;J}
\end{mbplan}

A feeder \textbf{splits nothing} --- it carries the whole bundle. A breakout is a thing you should be
able to see, and a patched box wearing coloured plates is exactly that; hiding one would leave a base
whose wiring could only be found by digging.

\begin{ietip}
A feeder is scenery, not hardware. It is never a node in the wire graph, holds no energy and costs
nothing per tick --- a run crossing four of them is still one connection between the same two boxes,
four blocks longer.
\end{ietip}

\section{Redstone channels}

A face set to \textbf{read} puts the signal arriving at it onto its conductor; a face set to
\textbf{emit} puts its conductor's signal back out. The signal reaches every box on the run, so a
lever at one end of a corridor throws a lamp at the other along the same bundle already carrying the
lighting.

A face does one of the three and never two. Reading and emitting on one face is how a redstone
network latches itself on and never lets go.

\section{Capacity and cost}

Each conductor carries what a conductor carries. A bundle is not a shared pipe with an allowance to
divide up: sixteen circuits down one corridor get sixteen circuits' worth of throughput.

Sixteen conductors are \emph{one} connection in the wire graph rather than sixteen, and a junction
box carrying nothing costs one integer comparison per tick. Replacing sixteen catenaries with one
conduit costs nothing but the crafting.

Energy moves as a bucket brigade: each box hands half the difference to whichever neighbour holds
less, and drains in full into any connector on the matching breakout. The honest cost of that is one
tick per hop, and boxes exist only at corners and ends.
